\subsection{Extension from a two-plane}
\label{sec:gct-extension}

Let \(V\) be the nine-dimensional symmetric tangent space at \(J_4\).
A three-plane \(L\) intersects any hyperplane \(H\subset V\) in
dimension at least two. Thus the search for \(L\) can start with
level two-planes \(W\subset H\). Fix a basis \(A,B\) of \(W\) and
let
\begin{equation*}
  E_W(C)=
  \left([z]q_j(sA+tB+zC)\right)_{j=2,3,4}.
\end{equation*}
This is a linear map with nine scalar rows, two from the quadratic,
three from the cubic, and four from the quartic. It kills \(W\).
If \(W\subset L\) for a level three-plane, then \(E_W\) kills all of
\(L\), so \(\operatorname{rank}E_W\le6\).
For a polynomial family on a coordinate space over an infinite
field, Lean verifies linearity of the first coefficient, its vanishing
along a containing level space, and the rank-nullity bound.

The quadratic rows cut out \(W^\perp\) for the nondegenerate Hessian.
Equivalently the cubic and quartic rows act on the five-dimensional
space \(W^\perp/W\), with seven rows and rank at most four needed for
extension. A full-rank map there excludes an extension immediately.
Rank four leaves one possible projective extension direction.

More generally put \(K_W=\ker E_W\) and choose a complement
\(K_W=W\oplus C_W\). A three-plane through \(W\) is level precisely
when one can choose a nonzero \(C\in C_W\) such that
\begin{equation*}
  [s^a t^b z^c]q_j(sA+tB+zC)=0
  \qquad(j=2,3,4,\ c\ge2,\ a+b+c=j).
\end{equation*}
There are ten remaining coefficients and each is homogeneous in
the coordinates of \(C\). The condition is independent of the chosen
complement, since adding an element of \(W\) to \(C\) changes the pencil
by a linear substitution in \(s,t\). The known zero-diagonal
two-planes demonstrate the role of these higher coefficients. Their
first-order maps have rank below seven, and the higher equations
exclude every extension.

There is also an obstruction using only the first-order map. On the
projective locus of level two-planes let \(\mathcal D\) be the closed
locus defined by \(\operatorname{rank}E_W\le6\). The rank condition
is given by minors on Grassmannian charts and is independent of the
choice of frame. Any level three-plane would place the entire
\(\operatorname{Gr}(2,L)\cong\mathbb P^2\) inside \(\mathcal D\).
Consequently an upper bound of one on the geometric dimension of
\(\mathcal D\) excludes every complex three-plane. The same bound on
a geometric special fiber of the projective model over
\(\mathbb Z_{(p)}\) suffices by upper semicontinuity of fiber
dimension~\cite{stacksfiberdimension}.

\paragraph{Two-planes defined over the prime field.}
Over \(\mathbb F_5\) the executable normalizes each projective tangent
direction by its first nonzero coordinate and enumerates 3918 level
directions. There are 2340 unordered pairs of distinct directions
spanning level two-planes. To check a pair, the mixed quadratic
coefficient must vanish. The cubic and quartic already vanish on
the two coordinate directions. On a pencil with relative scalar
\(t\), the remaining terms divided by \(t\) have degree at most two.
Evaluating at \(t=1,2,3\) therefore checks all their coefficients.
Row reduction on each pair gives 156 distinct two-planes.

Of these, 44 have \(K_W/W\) of dimension zero, 96 have dimension
one, and 16 have dimension two. In each one-dimensional case a
nonzero homogeneous coefficient excludes the unique projective
point. In each two-dimensional case the residual coefficients
have greatest common divisor one on the affine coordinate chart
and a nonzero value in the remaining direction. Hence none of the
156 planes extends over \(\overline{\mathbb F}_5\).

This also excludes any complex level three-plane containing a
rational level two-plane. For each fixed rational \(W\), extend it
over \(\mathbb Z_{(5)}\) by properness of its Grassmannian.
The incidence of level three-planes containing that integral
two-plane is projective. Its geometric special fiber is empty by
the checker, so its complex fiber is empty. In particular a complex
candidate can contain at most one independent rational direction.

\subsection{Two Pl\"ucker cuts and the full extension equations}
\label{sec:gct-plucker}

For a two-plane with basis \(A,B\), write
\(p_{ij,kl}=A_{ij}B_{kl}-A_{kl}B_{ij}\). Consider the two linear
forms in Pl\"ucker coordinates
\begin{align*}
 \alpha&=-p_{12,13}+p_{01,23}+p_{13,23},\\
 \beta&=-p_{12,23}-p_{03,12}+p_{01,02}.
\end{align*}
In any proposed level three-plane \(L\), the two-planes form the full
projective plane \(\mathbb P(\bigwedge^2L)\) in Pl\"ucker coordinates.
Thus at least one of them satisfies both cuts.

The 112 rank-deficient two-planes defined over \(\mathbb F_5\) contain
no point satisfying both cuts. There is such a point over
\(\mathbb F_{25}=\mathbb F_5[w]/(w^2-2)\), with basis
\begin{equation*}
 A=\begin{pmatrix}1&0&1&2\\0&2&0&3\\1&0&1&2\\2&3&2&0\end{pmatrix},
 \qquad
 B=\begin{pmatrix}4+w&0&1&0\\0&0&0&0\\1&0&4-w&0\\0&0&0&0\end{pmatrix}.
\end{equation*}
The checker verifies both cuts and the permanent pencil identity.
It computes a six-row minor of the first-variation map with
determinant \(2w-1\ne0\). Its kernel contains \(A,B\) and the
independent direction
\begin{equation*}
 C=\begin{pmatrix}-2w-1&1&0&0\\1&4&1&4\\0&1&2w+4&0\\0&4&0&1\end{pmatrix}.
\end{equation*}
Thus the rank is exactly six, and \(C\) represents the unique
projective extension direction. The full pencil difference is
\(-s^2z^2+sz^2-2z^4+z^3-z^2\), so the quadratic coefficient
already excludes this direction. The example shows that the
first-order equations on the cut cannot certify its emptiness.
The higher extension coefficients remove this point.

\paragraph{Excluding a direction supported off the diagonal on one edge.}
For a symmetric block-diagonal matrix with two \(2\times2\)
blocks \(U,V\), write \(\ell,-\ell\) for their entry sums and \(p,q\)
for their permanents. In the pencil at \(J_4\), the terms with no
crossing between the blocks give the product of the two order-two
pencils. The terms with one crossing in each direction give
\((4+t\ell)(4-t\ell)\), and the terms with two crossings in each
direction give four. The positive-degree coefficients therefore are
\begin{equation*}
 q_2=2(p+q-\ell^2),\qquad q_3=\ell(q-p),\qquad q_4=pq.
\end{equation*}
Their simultaneous vanishing forces \(p=q=\ell=0\). For one
nonzero off-diagonal pair, scale its entry to one and relabel the
pair as \((1,3)\). The other block is diagonal with zero sum and
zero product, hence zero. The active diagonal entries are
\(-1+h,-1-h\) with \(h^2=2\).

Projection onto the six off-diagonal entries is injective on any
level subspace in characteristic zero or five. Indeed a diagonal
direction has coefficients \((4-j)!e_j\) in terms of the elementary
symmetric functions of the diagonal entries. If they all vanish,
the entries are all zero. If a three-plane contains the normalized
one-edge direction \(C_h\), choose a basis \(A,B\) of its kernel
under the \((1,3)\) entry. Their projection onto the five remaining
off-diagonal slots is injective, so some pair of those slots is
a pivot. Relabelings preserving the edge, with \(h\mapsto-h\) when
its endpoints swap, reduce the ten pairs to
\[
 ((01),(02)),\quad ((01),(03)),\quad ((01),(12)),\quad ((01),(23)).
\]
The quadratic polar equations with \(C_h\) eliminate the \((3,3)\)
entry of each direction using the unit \(1-h\). The four remaining
systems include \(h^2=2\) and all pencil coefficients. Exact
Gr\"obner bases modulo 5 are \([1]\) in all four charts. Thus no
geometric three-plane at 5 contains such a direction. The incidence
of three-planes containing \(C_h\) is projective over the finite
base defined by \(h^2=2\) over \(\mathbb Z_{(5)}\). Proper
specialization gives the same exclusion over \(\mathbb C\).

\paragraph{The remaining charts.}
Write the six off-diagonal coordinates in row order as
\((x,y,z,b,c,e)=(P,R,v+S,Q,u-P,v-Q)\). In these coordinates the
alternating forms are \(u\wedge v-P\wedge Q\) and
\(P\wedge R+Q\wedge S\). By the single-edge exclusion, projection
of a level three-plane to \((u,v,P,Q,S)\) is injective. The kernel
of the projection consists of diagonal directions and directions
with only the \((0,2)\) off-diagonal entry available.

The projection of a cut two-plane is a two-plane in these five
coordinates that is isotropic for \(u\wedge v-P\wedge Q\). Scale
the five coordinates with respective weights \((3,-2,0,1,6)\).
The two terms of the alternating form have the same weight.
The lowest-weight projective limit therefore remains isotropic
and is fixed by the scaling, so it is a coordinate two-plane.
Its pivot pair cannot be \((u,v)\) or \((P,Q)\), and its nonzero
minor was nonzero on the original projected plane. This gives
eight possible pairs. Normalize the two pivot rows of \(A,B\),
set those entries of \(C\) to zero modulo their span, and
normalize one of the three remaining projected entries of
\(C\). No candidate extension vanishes at all three, by injectivity.
The resulting 24 affine systems impose both cuts and all
permanent coefficients of the full three-variable pencil.
A geometric point in any system gives a level three-plane at 5,
and every such plane supplies a geometric point in some system.
Unit ideals in all 24 systems would settle nonexistence over
\(\mathbb C\) by proper specialization of the unrestricted
three-plane locus. Those systems remain to be decided.

\subsection{A single path chart}
\label{sec:gct-path}

The off-diagonal projection gives a smaller cover without a choice of
two-plane. It is injective on every level three-plane. Hence some
three distinct off-diagonal coordinate functionals give an invertible
projection minor on the plane. Relabeling the four indices acts on
the 20 triples of edges with three orbits, represented by a star, a
triangle, and a path. The orbit sizes are four, four, and twelve.

Suppose no path minor is nonzero. Choose a nonzero star or triangle
minor and normalize its three rows to the identity. For each edge
outside the pivot triple, its coefficient on a given pivot direction
is, up to sign, the minor obtained by replacing that pivot with the
outside edge. The replacement is a path whenever the two edges meet.
Thus each outside edge is supported only on the direction of its
opposite pivot edge. Each of the three directions has both of its
opposite edges nonzero by the single-edge exclusion.

Each such direction therefore has two disjoint \(2\times2\) blocks.
The block calculation in Section~\ref{sec:gct-plucker} forces each
block with edge entry \(r\ne0\) to have diagonal
\(r(-1+h),r(-1-h)\), up to exchanging the entries, with \(h^2=2\).
Take any two of the directions and relabel their distinct matchings
as \((01),(23)\) and \((02),(13)\). Scale the first edge entry in
each to one and let \(x,y\) be their other edge entries. For each of
the sixteen sign choices of the four blocks, form the mixed
coefficients of the permanent at \(J_4\) in these two directions.
Together with \(h^2-2\) and \(rxy-1\), with \(r\) an auxiliary
variable, these coefficients generate the unit ideal. Over
\(\mathbb Q\) this is checked by exact Gr\"obner bases. Over
\(\mathbb F_5\) it is checked directly by row operations and
multiplication of the input polynomials through degree seven.
This contradicts the existence of the two directions.

Consequently every complex or \(\overline{\mathbb F}_5\) level
three-plane has a path minor. Relabel it to the edges \(01,12,23\),
and choose the frame whose projection onto these three entries is
the identity. The remaining system has 18 variables and 31
coefficients of degrees two through four. The executable emits this
single chart over \(\mathbb F_5\). A unit ideal would settle
nonexistence over \(\mathbb C\) by proper specialization. The
unrestricted path system remains open.
